% Source: Weinberg GR, physical PDF pp. 151--152
%         (printed pp. 126--127).
% Coverage: Section 5.3, Energy-Momentum Tensor; equations (5.3.1)--(5.3.9).
% Figures/tables/footnotes: none.
% Status: source-reviewed and compile-clean.
% Uncertainties: none.

\subsection{Energy-Momentum Tensor}\label{sec:5.3}

The density and current of energy and momentum were united in
Section~\ref{sec:2.8} into a symmetric tensor \(T^{\mu\nu}\) satisfying the
conservation equations
\begin{equation}
  \partial_\mu T^{\mu\nu}=G^\nu,
  \tag{5.3.1}\label{eq:5.3.1}
\end{equation}
where \(G^\nu\) is the density of the external force \(f^\nu\) acting on the
system. (For an isolated system, \(G^\nu=0\).) Define \(T^{\mu\nu}\) and
\(G^\nu\) as contravariant tensors that reduce to the special-relativistic
\(T^{\mu\nu}\) and \(G^\nu\) in the absence of gravitation. Then the
generally covariant equation that agrees with Eq.~\eqref{eq:5.3.1} in
locally inertial systems is
\begin{equation}
  \nabla_\mu T^{\mu\nu}=G^\nu,
  \tag{5.3.2}\label{eq:5.3.2}
\end{equation}
or, using Eq.~\eqref{eq:4.7.9},
\begin{equation}
  \frac{1}{\sqrt{-g}}\,
  \partial_\mu\!\left(\sqrt{-g}\,T^{\mu\nu}\right)
  =
  G^\nu
  -
  \Gamma^\nu{}_{\mu\lambda}T^{\mu\lambda}.
  \tag{5.3.3}\label{eq:5.3.3}
\end{equation}
The factor \(\sqrt{-g}\) is familiar from electrodynamics, and arises from
the fact that the invariant volume is \(\sqrt{-g}\,\dd^4x\). In contrast,
the second term on the right represents a gravitational force density. Just
as we would expect, this force depends on the system on which it acts only
through the energy-momentum tensor.

For a system of point particles the special-relativistic energy-momentum
tensor is given in Section~\ref{sec:2.8} as
\begin{equation}
  T^{\mu\nu}(x)
  =
  \sum_n m_n
  \int
  \frac{\dd x_n^\mu}{\dd\tau}\,
  \dd x_n^\nu\,
  \delta^4\!\left(x-x_n\right),
  \tag{5.3.4}\label{eq:5.3.4}
\end{equation}
the integral again being taken along the particle trajectory. Following
precisely the same reasoning as we did for \(J^\mu\) in the last section, we
conclude that the contravariant tensor that agrees with
Eq.~\eqref{eq:5.3.4} in the absence of gravitation is
\begin{equation}
  T^{\mu\nu}(x)
  =
  \frac{1}{\sqrt{-g(x)}}
  \sum_n m_n
  \int
  \frac{\dd x_n^\mu}{\dd\tau}\,
  \dd x_n^\nu\,
  \delta^4\!\left(x-x_n\right).
  \tag{5.3.5}\label{eq:5.3.5}
\end{equation}

For an electromagnetic field \(F^{\mu\nu}\) the special-relativistic
energy-momentum tensor was calculated in Section~\ref{sec:2.8} as
\begin{equation}
  T^{\mu\nu}_{\mathrm{em}}
  =
  F^\mu{}_\rho F^{\nu\rho}
  -
  \frac14\eta^{\mu\nu}F_{\rho\sigma}F^{\rho\sigma}.
  \tag{5.3.6}\label{eq:5.3.6}
\end{equation}
It takes no effort to see that the contravariant tensor that agrees with
Eq.~\eqref{eq:5.3.6} in the absence of gravitation is
\begin{equation}
  T^{\mu\nu}_{\mathrm{em}}
  =
  F^\mu{}_\rho F^{\nu\rho}
  -
  \frac14 g^{\mu\nu}F_{\rho\sigma}F^{\rho\sigma}.
  \tag{5.3.7}\label{eq:5.3.7}
\end{equation}
For a system consisting of particles and radiation, the energy-momentum
tensor is the sum of Eqs.~\eqref{eq:5.3.5} and \eqref{eq:5.3.7}.

Returning for a moment to the purely material energy-momentum tensor
\eqref{eq:5.3.5}, we easily compute that
\[
  \int
  T^{\mu0}\sqrt{-g}\,\dd^3x
  =
  \sum_n m_n\frac{\dd x_n^\mu}{\dd\tau},
\]
the sum running over all particles in the volume of integration. This
suggests that \(T^{\mu0}\sqrt{-g}\) is to be regarded in general as the
spatial density of energy and momentum. In particular, we are tempted to
define the energy, momentum, and angular momentum for an arbitrary system by
\begin{equation}
  P^\mu
  \equiv
  \int
  T^{\mu0}\sqrt{-g}\,\dd^3x,
  \tag{5.3.8}\label{eq:5.3.8}
\end{equation}
\begin{equation}
  J^{\mu\nu}
  \equiv
  \int
  \left(
    x^\mu T^{\nu0}
    -
    x^\nu T^{\mu0}
  \right)
  \sqrt{-g}\,\dd^3x.
  \tag{5.3.9}\label{eq:5.3.9}
\end{equation}
However, these quantities are \emph{not} contravariant tensors and are
\emph{not} conserved, because \(T^{\mu\nu}\sqrt{-g}\) is not conserved;
that is, \(\partial_\nu(\sqrt{-g}\,T^{\mu\nu})\) does not vanish, owing to
the exchange of energy and momentum between matter and gravitation.
